% Spellchecking done with grammarly
\section{Baseline Methodology}
The purpose of the underlying paper was to create an algorithm which solves the issue of plotting vibration data on a line graph where the observer can't distinguish between noise and the actual vibrations. The overall goal was to: further enhance and complement the visual inspection in the time
domain, enabling users to quickly gain an overview of different
vibration profiles within the signal or to detect periodicity within
noise. (\cite{rakuschek2025visual} p.1). In the background this was done by using  time delay embeddings, which is a method that takes the time series data and transforms it into a point cloud which can than be projected onto a 2D plane. Another really important aspect of the algorithm which is to be transformed is the use of principal component analysis. It is a dimensionality reduction technique that helps reduce the number of features in a dataset while keeping the most important information. As stated by \textcite{KarlPearson}, principal component analysis simplifies complex datasets by transforming correlated features into a smaller set of uncorrelated components. When the two algorithms are applied on a harmonic oscillation the result is a ellipsoid and as more noise is added its really hard to distinguished between noise or the hidden signal. However, applying time delay embedding with appropriate window sizes effectively reveals the underlying vibration signal.